% !TeX program = xelatex
% Kompilasi dari folder proyek ini: .\Compile-Tugas.ps1
\documentclass[10pt]{article}
% !TeX root = review-jurnal.tex
\newif\ifTampilanWord
% true: mengikuti ukuran halaman, margin, dan ukuran huruf aktual Word asli.
% false: A4, margin 2 cm, semua huruf minimal 10 pt (ketentuan pengumpulan).
% Peralihan ke false dapat mengubah jumlah halaman; periksa setelah diisi.
\TampilanWordfalse

\usepackage{geometry}
\ifTampilanWord
  % Ukuran dan margin aktual DOCX (twips dikonversi ke bp).
  \geometry{paperwidth=612bp,paperheight=792bp,left=56.7bp,right=56.7bp,top=48.2bp,bottom=45.35bp,footskip=17bp}
\else
  \geometry{a4paper,margin=2cm,footskip=20bp}
\fi
\usepackage{fontspec}
\setmainfont{Times New Roman}
\usepackage{tabularx,array,xfp,fancyhdr}
\usepackage{ragged2e}
\usepackage[table]{xcolor}
\usepackage[hidelinks]{hyperref}
\urlstyle{same}
\setlength{\parindent}{0pt}
\setlength{\parskip}{0pt}
\setlength{\emergencystretch}{2em}
\newcommand{\ukuranword}[1]{%
  \ifTampilanWord
    \fontsize{#1bp}{\fpeval{#1*1.0}bp}\selectfont
  \else
    \ifdim#1pt<10pt \fontsize{10bp}{10bp}\selectfont
    \else
      \ifdim#1pt>12pt \fontsize{12bp}{12bp}\selectfont
      \else \fontsize{#1bp}{\fpeval{#1*1.0}bp}\selectfont\fi
    \fi
  \fi}
\newcommand{\teksword}[3]{\par\vspace{#2bp}{\ukuranword{#1}#3\par}\vspace{0bp}}
\newcommand{\judulword}[3]{\par\addvspace{#2bp}{\ukuranword{#1}\bfseries #3\par}\nobreak\vspace{0bp}\nobreak}
\newcommand{\pertanyaan}[2]{\par{\ukuranword{#1}\leftskip=9.9bp\hangindent=-8.5bp #2\par}\nobreak}
% Biarkan kosong untuk mereproduksi lembar tugas. Isi dengan analisis sendiri.
\newcommand{\jawaban}[1]{\if\relax\detokenize{#1}\relax\else\par{\ukuranword{10}\justifying\setlength{\parindent}{0pt}\noindent #1\par}\vspace{0bp}\fi}
\newcommand{\garisisian}{\teksword{8.5}{0}{................................................................................................................................................}}
\newcolumntype{Y}{>{\raggedright\arraybackslash}X}
\definecolor{headerword}{HTML}{D9EAF7}
\setlength{\tabcolsep}{5.4bp}
\setlength{\arrayrulewidth}{0.5bp}
\renewcommand{\arraystretch}{1}
\NewDocumentEnvironment{tabelword}{m+b}{%
  \par\begingroup\ukuranword{7.5}%
  \setlength{\extrarowheight}{2bp}%
  \begin{tabularx}{\linewidth}{#1}\hline #2\end{tabularx}\par\endgroup\vspace{0bp}%
}{}
\pagestyle{fancy}
\fancyhf{}
\renewcommand{\headrulewidth}{0pt}
\renewcommand{\footrulewidth}{0pt}
\fancyfoot[C]{{\ukuranword{8}Review Dua Artikel Ilmiah | Halaman \thepage}}
% Ganti teks isian dengan data Anda pada salinan pribadi.
% Pertahankan kurung kurawal { } sebagai sintaks LaTeX.
\newcommand{\NamaMahasiswa}{.........}
\newcommand{\NIMMahasiswa}{.........}
\newcommand{\KelasMahasiswa}{.........}
\newcommand{\DosenPengampu}{.........}
\newcommand{\ProgramStudi}{.........}
\newcommand{\MataKuliah}{.........}

\begin{document}
\raggedright
\teksword{15}{0}{\centering\bfseries REVIEW JURNAL}
\teksword{11}{0}{\centering\bfseries ANALISIS DAN PERBANDINGAN DUA ARTIKEL ILMIAH}
\teksword{9}{0}{Program Studi: \ProgramStudi\quad Mata Kuliah: \MataKuliah}
\teksword{9}{0}{Nama: \NamaMahasiswa\quad NIM: \NIMMahasiswa\quad Kelas: \KelasMahasiswa}
\teksword{9}{0}{Dosen Pengampu: \DosenPengampu}
\judulword{10}{4}{PANDUAN REVIEW}
\pertanyaan{9}{1. Tinjau dua artikel ilmiah sesuai bidang atau topik yang dipilih. Topik: ......... Rentang tahun: .........}
\pertanyaan{9}{2. Kedua artikel harus memiliki keterkaitan yang memungkinkan perbandingan masalah, metode, atau temuan.}
\pertanyaan{9}{3. Analisis setiap jurnal secara terpisah, kemudian bandingkan keduanya dan buat usulan pengembangan.}
\pertanyaan{9}{4. Gunakan sumber ilmiah yang dapat dipertanggungjawabkan. Dilarang hanya menyalin abstrak atau kesimpulan.}
\pertanyaan{9}{5. Lampirkan file atau tautan kedua artikel jurnal. Simpan laporan dalam format PDF.}

\judulword{10}{4}{A. REVIEW JURNAL 1}
\teksword{8.5}{0}{Judul: ..........}
\teksword{8.5}{0}{Penulis: ......... Tahun: ......... Jurnal: ..........}
\teksword{8.5}{0}{Volume/nomor/halaman: ......... DOI/URL: ..........}

\judulword{10}{4}{1. Masalah dan Tujuan Penelitian}
\pertanyaan{9}{1. Apa masalah utama yang ingin diselesaikan dan mengapa penting diteliti?}
\teksword{8.5}{0}{.........}
\pertanyaan{9}{2. Apa keterbatasan penelitian sebelumnya serta dampaknya?}
\teksword{8.5}{0}{.........}
\pertanyaan{9}{3. Apa tujuan khusus penelitian?}
\teksword{8.5}{0}{.........}

\judulword{10}{4}{2. Metode Penelitian}
\pertanyaan{9}{1. Pendekatan dan metode apa yang digunakan serta mengapa dipilih peneliti?}
\teksword{8.5}{0}{.........}
\pertanyaan{9}{2. Uraikan desain penelitian, tahapan pelaksanaan, serta objek atau subjek yang diteliti.}
\teksword{8.5}{0}{.........}
\pertanyaan{9}{3. Bagaimana data dikumpulkan dan dianalisis untuk menjawab pertanyaan penelitian?}
\teksword{8.5}{0}{.........}
\pertanyaan{9}{4. Apakah metode sesuai dengan masalah? Jelaskan kelebihan dan keterbatasannya.}
\teksword{8.5}{0}{.........}

\judulword{10}{4}{3. Data, Bahan, dan Perangkat Penelitian}
\begin{tabelword}{|Y|Y|Y|}
\rowcolor{headerword}
\textbf{\ukuranword{8}Aspek penelitian} & \textbf{\ukuranword{8}Rincian yang digunakan} & \textbf{\ukuranword{8}Peran dalam penelitian} \\ \hline
Data/sumber data & ......... & ......... \\ \hline
Sampel/objek/subjek & ......... & ......... \\ \hline
Bahan/instrumen & ......... & ......... \\ \hline
Perangkat/alat analisis & ......... & ......... \\ \hline
\end{tabelword}

\judulword{10}{4}{4. Hasil dan Evaluasi}
\pertanyaan{9}{1. Bagaimana hasil dievaluasi dan indikator apa yang digunakan?}
\teksword{8.5}{0}{.........}
\pertanyaan{9}{2. Tuliskan temuan utama beserta data, ukuran hasil, atau bukti yang dilaporkan.}
\teksword{8.5}{0}{.........}
\pertanyaan{9}{3. Apakah temuan mendukung tujuan atau kesimpulan penelitian? Jelaskan bukti dan batasannya.}
\teksword{8.5}{0}{.........}

\judulword{10}{4}{5. Kelebihan, Kekurangan, dan Pengembangan}
\teksword{8.5}{0}{.........}
\teksword{8.5}{0}{.........}
\teksword{8.5}{0}{.........}


\newpage
\judulword{12}{0}{B. REVIEW JURNAL 2}
\teksword{8.5}{0}{Judul: ..........}
\teksword{8.5}{0}{Penulis: ......... Tahun: ......... Jurnal: ..........}
\teksword{8.5}{0}{Volume/nomor/halaman: ......... DOI/URL: ..........}

\judulword{10}{4}{1. Masalah dan Tujuan Penelitian}
\pertanyaan{9}{1. Apa masalah utama yang ingin diselesaikan dan mengapa penting diteliti?}
\teksword{8.5}{0}{.........}
\pertanyaan{9}{2. Apa keterbatasan penelitian sebelumnya serta dampaknya?}
\teksword{8.5}{0}{.........}
\pertanyaan{9}{3. Apa tujuan khusus penelitian?}
\teksword{8.5}{0}{.........}

\judulword{10}{4}{2. Metode Penelitian}
\pertanyaan{9}{1. Pendekatan dan metode apa yang digunakan serta mengapa dipilih peneliti?}
\teksword{8.5}{0}{.........}
\pertanyaan{9}{2. Uraikan desain penelitian, tahapan pelaksanaan, serta objek atau subjek yang diteliti.}
\teksword{8.5}{0}{.........}
\pertanyaan{9}{3. Bagaimana data dikumpulkan dan dianalisis untuk menjawab pertanyaan penelitian?}
\teksword{8.5}{0}{.........}
\pertanyaan{9}{4. Apakah metode sesuai dengan masalah? Jelaskan kelebihan dan keterbatasannya.}
\teksword{8.5}{0}{.........}

\judulword{10}{4}{3. Data, Bahan, dan Perangkat Penelitian}
\begin{tabelword}{|Y|Y|Y|}
\rowcolor{headerword}
\textbf{\ukuranword{8}Aspek penelitian} & \textbf{\ukuranword{8}Rincian yang digunakan} & \textbf{\ukuranword{8}Peran dalam penelitian} \\ \hline
Data/sumber data & ......... & ......... \\ \hline
Sampel/objek/subjek & ......... & ......... \\ \hline
Bahan/instrumen & ......... & ......... \\ \hline
Perangkat/alat analisis & ......... & ......... \\ \hline
\end{tabelword}

\judulword{10}{4}{4. Hasil dan Evaluasi}
\pertanyaan{9}{1. Bagaimana hasil dievaluasi dan indikator apa yang digunakan?}
\teksword{8.5}{0}{.........}
\pertanyaan{9}{2. Tuliskan temuan utama beserta data, ukuran hasil, atau bukti yang dilaporkan.}
\teksword{8.5}{0}{.........}
\teksword{8.5}{0}{.........}
\pertanyaan{9}{3. Apakah temuan mendukung tujuan atau kesimpulan penelitian? Jelaskan bukti dan batasannya.}
\teksword{8.5}{0}{.........}

\judulword{10}{4}{5. Kelebihan, Kekurangan, dan Pengembangan}
\teksword{8.5}{0}{.........}
\teksword{8.5}{0}{.........}
\teksword{8.5}{0}{.........}


\newpage
\judulword{11}{0}{C. PERBANDINGAN DAN ANALISIS KRITIS KEDUA JURNAL}
\teksword{8.5}{0}{Bandingkan kedua artikel. Jika aspek tidak dijelaskan, tuliskan “tidak dijelaskan”.}
\begin{tabelword}{|Y|Y|Y|}
\rowcolor{headerword}
\textbf{\ukuranword{8}Aspek} & \textbf{\ukuranword{8}Jurnal 1} & \textbf{\ukuranword{8}Jurnal 2} \\ \hline
Masalah dan tujuan & ......... & ......... \\ \hline
Metode penelitian & ......... & ......... \\ \hline
Teknik analisis & ......... & ......... \\ \hline
Data dan sumber & ......... & ......... \\ \hline
Sampel/objek/subjek & ......... & ......... \\ \hline
Instrumen/perangkat & ......... & ......... \\ \hline
Desain evaluasi & ......... & ......... \\ \hline
Indikator evaluasi & ......... & ......... \\ \hline
Hasil utama & ......... & ......... \\ \hline
Kelebihan & ......... & ......... \\ \hline
Keterbatasan & ......... & ......... \\ \hline
\end{tabelword}

\judulword{10}{4}{Pertanyaan Analisis Perbandingan}
\pertanyaan{9}{1. Apa persamaan dan perbedaan utama kedua penelitian dari sisi masalah, metode, dan bukti?}
\teksword{8.5}{0}{.........}
\pertanyaan{9}{2. Metode mana yang paling sesuai untuk masalah yang diteliti? Jelaskan berdasarkan isi jurnal.}
\teksword{8.5}{0}{.........}
\pertanyaan{9}{3. Penelitian mana yang temuannya lebih meyakinkan? Pertimbangkan kualitas data, metode, dan bukti.}
\teksword{8.5}{0}{.........}
\pertanyaan{9}{4. Apakah kedua pendekatan dapat digabungkan atau saling melengkapi? Jelaskan.}
\teksword{8.5}{0}{.........}

\judulword{10}{4}{D. RESEARCH GAP}
\pertanyaan{9}{1. Apa masalah atau keterbatasan yang belum diselesaikan oleh kedua penelitian?}
\teksword{8.5}{0}{.........}
\pertanyaan{9}{2. Aspek apa yang masih dapat ditingkatkan dari desain, data, metode analisis, atau evaluasi?}
\teksword{8.5}{0}{.........}
\pertanyaan{9}{3. Metode atau pendekatan alternatif apa yang berpotensi memperbaiki hasil?}
\teksword{8.5}{0}{.........}



\newpage
\judulword{11}{0}{E. USULAN PENGEMBANGAN DAN KESIMPULAN}
\judulword{10}{4}{Usulan Pengembangan Penelitian}
\pertanyaan{9}{1. Judul/tema pengembangan yang diusulkan:}
\teksword{8.5}{0}{.........}
\pertanyaan{9}{2. Permasalahan yang akan diselesaikan:}
\teksword{8.5}{0}{.........}
\pertanyaan{9}{3. Metode penelitian yang dipilih dan alasan pemilihannya:}
\teksword{8.5}{0}{.........}
\pertanyaan{9}{4. Data, bahan, dan perangkat yang dibutuhkan:}
\teksword{8.5}{0}{.........}
\pertanyaan{9}{5. Gambaran singkat pelaksanaan penelitian yang diusulkan:}
\teksword{8.5}{0}{.........}
\pertanyaan{9}{6. Rencana evaluasi dan indikator keberhasilan:}
\teksword{8.5}{0}{.........}
\teksword{8.5}{0}{.........}

\judulword{10}{4}{Kesimpulan Akhir}
\teksword{8.5}{0}{Rangkum hasil review kedua jurnal: masalah, metode, hasil, kelebihan-keterbatasan, dan alasan usulan pengembangan.}
\teksword{8.5}{0}{.........}
\teksword{8.5}{0}{.........}

\judulword{10}{4}{Daftar Pustaka}
\teksword{8.5}{0}{Cantumkan kedua artikel utama dan referensi pendukung dengan gaya IEEE atau APA secara konsisten.}
\teksword{10}{0}{[1] ..........}
\teksword{10}{0}{[2] ..........}

\judulword{10}{4}{Contoh Format dan Penggunaan}
\pertanyaan{9}{1. Kertas A4; Times New Roman 10–12 pt; spasi 1,0–1,15; margin 2 cm.}
\pertanyaan{9}{2. Contoh panjang laporan: 4 halaman isi; sesuaikan dengan kebutuhan atau ketentuan institusi.}
\pertanyaan{9}{3. Nama file: ReviewJurnal\_NIM\_Nama.pdf.}
\pertanyaan{9}{4. Presentasikan hasil review jika diperlukan sesuai konteks penggunaan.}
\judulword{10}{4}{CONTOH RUBRIK PENILAIAN}
\teksword{8.5}{0}{Komponen dan bobot berikut adalah contoh; sesuaikan dengan kebutuhan, dengan total bobot 100\%.}
\begin{tabelword}{|Y|Y|}
\rowcolor{headerword}
\textbf{\ukuranword{8}Komponen} & \textbf{\ukuranword{8}Bobot} \\ \hline
Analisis masalah kedua jurnal & 15\% \\ \hline
Analisis metode penelitian & 20\% \\ \hline
Analisis data, bahan, dan perangkat & 10\% \\ \hline
Analisis hasil dan bukti & 15\% \\ \hline
Kelebihan dan keterbatasan & 10\% \\ \hline
Perbandingan kritis kedua jurnal & 15\% \\ \hline
Research gap dan usulan pengembangan & 10\% \\ \hline
Kesimpulan, referensi, dan tata tulis & 5\% \\ \hline
TOTAL & 100\% \\ \hline
\end{tabelword}
\teksword{9}{0}{Dosen Pengampu,}
\teksword{9}{24}{(\DosenPengampu)}
\teksword{9}{0}{NIDN: .........}

\end{document}
